\documentclass[aspectratio=169, xcolor=dvipsnames]{beamer}

% --- Theme Setup ---
\usetheme{Madrid}
\usecolortheme{default}

% --- Packages ---
\usepackage{graphicx}
\usepackage{xcolor}
\usepackage{booktabs}
\usepackage{hyperref}
\usepackage{adjustbox}
\usepackage{amssymb}
\usepackage{amsmath}
\usepackage{listings}
\usepackage{caption}
\usepackage{tikz}

\lstset{
  basicstyle=\ttfamily\small,
  keywordstyle=\color{blue}\bfseries,
  stringstyle=\color{red},
  showstringspaces=false,
  breaklines=true
}

% --- Title Information ---
\title{Module 35}
\subtitle{Application Design and Development/5: Application Development and Mobile}
\author{Partha Pratim Das}
\institute{Department of Computer Science and Engineering\\Indian Institute of Technology, Kharagpur\\ppd@cse.iitkgp.ac.in}
\date{\today}

\begin{document}

% Title Slide
\begin{frame}
    \titlepage
\end{frame}

\begin{frame}{Module Recap}
    \begin{itemize}
        \item Learnt how to accessing PostgreSQL database from Python
        \item Learnt to build Python Web Application with PostgreSQL and Flask
    \end{itemize}
\end{frame}

\begin{frame}{Module Objectives}
    \begin{itemize}
        \item To explore the Rapid Application Development Process
        \item To understand the issues in Application Performance and Security
        \item To understand the similarity and differences between how Mobile Apps and Web Applications
    \end{itemize}
\end{frame}

\begin{frame}{Module Outline}
    \begin{itemize}
        \item Rapid Application Development
        \item Application Performance and Security
        \item Mobile Apps
    \end{itemize}
\end{frame}

\section{Rapid Application Development}
\begin{frame}
  \vfill
  \centering
  \begin{beamercolorbox}[sep=8pt,center,shadow=true,rounded=true]{title}
    \usebeamerfont{title}Rapid Application Development\par%
  \end{beamercolorbox}
  \vfill
\end{frame}

\begin{frame}{Rapid Application Development}
    \begin{itemize}
        \item A lot of effort is required to develop Web application interfaces, especially the rich interaction functionality associated with Web 2.0 applications
        \item Several approaches to speed up application development
        \begin{itemize}
            \item Function library to generate user-interface elements
            \item Drag-and-drop features in an IDE to create user-interface elements
            \item Automatically generate code for user interface from a declarative specification
        \end{itemize}
        \item Used as part of \textbf{Rapid Application Development (RAD)} tools even before Web
        \item RAD Software is an agile model that focuses on fast prototyping and quick feedback in app development to ensure speedier delivery and an efficient result
        \item App development has 4 phases: business modeling, data modeling, process modeling, and testing \& turnover: Defining the requirements, Prototyping, Receiving feedback and Finalizing the software
        \item With RAD, the time between prototypes and iterations is short, and integration occurs since inception.
    \end{itemize}
\end{frame}

\begin{frame}{Rapid Application Development (2)}
    \begin{itemize}
        \item \textbf{Web application development frameworks}
        \begin{itemize}
            \item \textbf{Java Server Faces (JSF)}
            \begin{itemize}
                \item $\triangleright$ A set of APIs for representing UI components and managing their state, handling events and input validation, defining page navigation, and supporting internationalization and accessibility
                \item $\triangleright$ JSP custom tag library for expressing a JSF interface within a JSP page
            \end{itemize}
            \item \textbf{Ruby on Rails}
            \begin{itemize}
                \item $\triangleright$ Allows easy creation of simple CRUD (create, read, update and delete) interfaces by code generation from database schema or object model
            \end{itemize}
        \end{itemize}
        \item \textbf{RAD Platforms and Tools}
        \begin{itemize}
            \item G Suite
            \item Google App Engine
            \item Microsoft Azure
            \item Amazon Elastic Compute Cloud (EC2)
            \item AWS Elastic Beanstalk
        \end{itemize}
    \end{itemize}
\end{frame}

\begin{frame}{ASP.NET and Visual Studio}
    \begin{itemize}
        \item ASP.NET provides a variety of controls that are interpreted at server, and generate HTML code
        \item Visual Studio provides drag-and-drop development using these controls
        \begin{itemize}
            \item For example, menus and list boxes can be associated with \texttt{DataSet} object
        \end{itemize}
        \item Validator controls (constraints) can be added to form input fields
        \begin{itemize}
            \item $\triangleright$ JavaScript to enforce constraints at client, and separately enforced at server
        \end{itemize}
        \item User actions such as selecting a value from a menu can be associated with actions at server
        \item DataGrid provides convenient way of displaying SQL query results in tabular format
    \end{itemize}
\end{frame}

\section{Application Performance and Security}
\begin{frame}
  \vfill
  \centering
  \begin{beamercolorbox}[sep=8pt,center,shadow=true,rounded=true]{title}
    \usebeamerfont{title}Application Performance and Security\par%
  \end{beamercolorbox}
  \vfill
\end{frame}

\begin{frame}{Application Performance}
    \begin{itemize}
        \item Performance is an issue for popular Web sites
        \begin{itemize}
            \item May be accessed by millions of users every day, thousands of requests per second at peak time
        \end{itemize}
        \item Caching techniques used to reduce cost of serving pages by exploiting commonalities between requests
        \item At the server site:
        \begin{itemize}
            \item $\triangleright$ Caching of JDBC connections between servlet requests
            \begin{itemize}
                \item a.k.a. connection pooling
            \end{itemize}
            \item $\triangleright$ Caching results of database queries
            \begin{itemize}
                \item Cached results must be updated if underlying database changes
            \end{itemize}
            \item $\triangleright$ Caching of generated HTML
        \end{itemize}
        \item At the client's network
        \begin{itemize}
            \item $\triangleright$ Caching of pages by Web proxy
        \end{itemize}
    \end{itemize}
\end{frame}

\begin{frame}[fragile]{Application Security: SQL Injection}
    \begin{itemize}
        \item Suppose query is constructed using
        \begin{center}\texttt{"select * from instructor where name = '" + name + "'"}\end{center}
        \item Suppose the user, instead of entering a name, enters:
        \begin{center}\texttt{X' or 'Y' = 'Y}\end{center}
        \item then the resulting statement becomes:
        \begin{center}\texttt{"select * from instructor where name = '" + "X' or 'Y' = 'Y" + "'"}\end{center}
        \item which is:
        \begin{itemize}
            \item $\triangleright$ \texttt{select * from instructor where name = 'X' or 'Y' = 'Y'}
        \end{itemize}
        \item User could have even used
        \begin{itemize}
            \item $\triangleright$ \texttt{X'; update instructor set salary = salary + 10000; -{}-}
        \end{itemize}
        \item Prepared statement internally uses:
        \begin{center}\texttt{select * from instructor where name = 'X \textbackslash' or \textbackslash'Y \textbackslash' = \textbackslash'Y'}\end{center}
        \item \textbf{Always use prepared statements, with user inputs as parameters}
        \item Is the following prepared statement secure?
        \begin{center}\texttt{conn.prepareStatement("select * from instructor where name = '" + name + "'")}\end{center}
    \end{itemize}
\end{frame}

\begin{frame}{Application Security (2): Password Leakage}
    \begin{itemize}
        \item Never store passwords, such as database passwords, in clear text in scripts that may be accessible to users
        \item For example, in files in a directory accessible to a web server
        \begin{itemize}
            \item $\triangleright$ Normally, web server will execute, but not provide source of script files such as \texttt{file.jsp} or \texttt{file.php}, but source of editor backup files such as \texttt{file.jsp\~{}}, or \texttt{.file.jsp.swp} may be served
        \end{itemize}
        \item Restrict access to database server from IPs of machines running application servers
        \begin{itemize}
            \item Most databases allow restriction of access by source IP address
        \end{itemize}
    \end{itemize}
\end{frame}

\begin{frame}{Application Security (3): Authentication}
    \begin{itemize}
        \item Single factor authentication such as passwords too risky for critical applications
        \begin{itemize}
            \item guessing of passwords, sniffing of packets if passwords are not encrypted
            \item passwords reused by user across sites
            \item spyware which captures password
        \end{itemize}
        \item Two-factor authentication
        \begin{itemize}
            \item For example, password plus one-time password sent by SMS
            \item For example, password plus one-time password devices
            \begin{itemize}
                \item $\triangleright$ device generates a new pseudo-random number every minute, and displays to user
                \item $\triangleright$ user enters the current number as password
                \item $\triangleright$ application server generates same sequence of pseudo-random numbers to check that the number is correct.
            \end{itemize}
        \end{itemize}
    \end{itemize}
\end{frame}

\begin{frame}[fragile]{Application Security (4): Application-Level Authorization}
    \begin{itemize}
        \item Current SQL standard does not allow fine-grained authorization such as 'students can see their own grades, but not other's grades'
        \item \textbf{Problem 1:} Database has no idea who are application users
        \item \textbf{Problem 2:} SQL authorization is at the level of tables, or columns of tables, but not to specific rows of a table
        \item One workaround: use views such as
        \begin{verbatim}
create view studentTakes as
select * from takes
where takes.ID = syscontext.user_id()
        \end{verbatim}
        \item where \texttt{syscontext.user\_id()} provides end user identity
        \begin{itemize}
            \item $\triangleright$ end user identity must be provided to the database by the application
        \end{itemize}
        \item Having multiple such views is cumbersome
    \end{itemize}
\end{frame}

\begin{frame}{Application Security (5): Application-Level Authorization}
    \begin{itemize}
        \item Currently, authorization is done entirely in application
        \begin{itemize}
            \item Entire application code has access to entire database
            \item large surface area, making protection harder
        \end{itemize}
        \item Alternative: fine-grained (row-level) authorization schemes
        \begin{itemize}
            \item extensions to SQL authorization proposed but not currently implemented
            \item Oracle Virtual Private Database (VPD) allows predicates to be added transparently to all SQL queries, to enforce fine-grained authorization
            \begin{itemize}
                \item $\triangleright$ For example, add \texttt{ID = sys\_context.user\_id()} to all queries on \texttt{student} relation if user is a student
            \end{itemize}
        \end{itemize}
    \end{itemize}
\end{frame}

\begin{frame}{Application Security (6): Audit Trails}
    \begin{itemize}
        \item Applications must log actions to an audit trail, to detect who carried out an update, or accessed some sensitive data
        \item Audit trails used after-the-fact to
        \begin{itemize}
            \item detect security breaches
            \item repair damage caused by security breach
            \item trace who carried out the breach
        \end{itemize}
        \item Audit trails needed at
        \begin{itemize}
            \item Database level, and at
            \item Application level
        \end{itemize}
    \end{itemize}
\end{frame}

\section{Challenges in Web Application Development}
\begin{frame}
  \vfill
  \centering
  \begin{beamercolorbox}[sep=8pt,center,shadow=true,rounded=true]{title}
    \usebeamerfont{title}Challenges in Web Application Development\par%
  \end{beamercolorbox}
  \vfill
\end{frame}

\begin{frame}{Challenges in Web Application Development}
    \begin{itemize}
        \item User Interface and User Experience
        \item Scalability
        \item Performance
        \item Knowledge of Framework and Platforms
        \item Security
    \end{itemize}
    \vspace{2em}
    \small Source: 5 Challenges in Web Application Development
\end{frame}

\section{Mobile Apps}
\begin{frame}
  \vfill
  \centering
  \begin{beamercolorbox}[sep=8pt,center,shadow=true,rounded=true]{title}
    \usebeamerfont{title}Mobile Apps\par%
  \end{beamercolorbox}
  \vfill
\end{frame}

\begin{frame}{What is a Mobile App?}
    \begin{itemize}
        \item A type of application software designed to run on a mobile device, such as a smartphone or tablet computer
        \item Developed specifically for use on small, wireless computing devices, such as smartphones and tablets
        \item Designed with consideration for the demands and constraints of the devices and also to take advantage of any specialized capabilities
        \begin{itemize}
            \item - Form Factor - influences display and navigation
            \item - Limited Memory
            \item - Limited Computing Power
            \item - Limited Power
            \item - Limited Bandwidth
        \end{itemize}
        \item \dots
        \begin{itemize}
            \item[+] Availability of sensors like accelerometer
            \item[+] Availability of touchscreen - Gesture-based Navigation
        \end{itemize}
    \end{itemize}
\end{frame}

\begin{frame}{Mobile Website vis-\`a-vis Mobile App}
    \begin{itemize}
        \item \textbf{Mobile Website}
        \begin{itemize}
            \item Similar to any other website in that it consists of browser-based HTML pages
            \item Can display text content, data, images and video
            \item Typically accessed over WiFi or 3G or 4G networks
            \item Designed for the smaller handheld display and touch-screen interface
            \item Can also access mobile-specific features such as click-to-call (to dial a phone number) or location-based mapping
        \end{itemize}
        \item \textbf{Mobile Apps}
        \begin{itemize}
            \item Actual applications that are downloaded and installed on the mobile device
            \item Users download apps from device-specific portals such as App Store, Google Play Store
            \item The app may
            \begin{itemize}
                \item $\triangleright$ pull content and data from the Internet, in similar fashion to a website, or
                \item $\triangleright$ download the content so that it can be accessed without an Internet connection
            \end{itemize}
        \end{itemize}
    \end{itemize}
    \small Source: \url{https://www.slideshare.net/hassandar18/architecture-of-mobile-software-applications?from_action=save}
\end{frame}

\begin{frame}{Architecture of Mobile App}
    \begin{itemize}
        \item Typically 3 tier
        \begin{itemize}
            \item Presentation
            \item Business
            \item Data
        \end{itemize}
        \item Data Layer is often split between:
        \begin{itemize}
            \item Local Data
            \item Remote Data
        \end{itemize}
        \item Needs customization for platform
        \begin{itemize}
            \item Android
            \item iOS
            \item Windows
        \end{itemize}
    \end{itemize}
\end{frame}

\begin{frame}{Types of Mobile Apps}
    \begin{itemize}
        \item \textbf{Native Apps}: Completely written in the native language of a platform
        \begin{itemize}
            \item iOS $\rightarrow$ Objective-C; Android $\rightarrow$ Java or C/C++
            \item Platform specific (heavily dependent on OS)
        \end{itemize}
        \item \textbf{Web Apps}: Run completely inside of a Web browser.
        \begin{itemize}
            \item Features interfaces built with HTML or CSS
            \item Powered via Web programming languages $\rightarrow$ Ruby on Rails, JavaScript, PHP, or Python
            \item Portable across any phone, tablet, or computer
        \end{itemize}
        \item \textbf{Hybrid Apps}: Combines attributes of both native and Web apps.
        \begin{itemize}
            \item Attempts to use redundant, common code that can be used across platforms, and
            \item Tailors required attributes to the native system
        \end{itemize}
    \end{itemize}
    \small Source: \url{https://www.slideshare.net/hassandar18/architecture-of-mobile-software-applications?from_action=save}
\end{frame}

\begin{frame}{Design Issues}
    \begin{itemize}
        \item Determine Device
        \item Note Device Resources - memory, power, speed
        \item Consider Bandwidth
        \item Decide on Architecture Layers
        \item Select Technology
        \item Define User Interface
        \item Select Navigation
        \item Maintain Flow
    \end{itemize}
    \vspace{2em}
    \small Source: \url{https://www.slideshare.net/hassandar18/architecture-of-mobile-software-applications?from_action=save}
\end{frame}

\begin{frame}{Module Summary}
    \begin{itemize}
        \item Understood the steps in the Rapid Application Development Process
        \item Exposed to the issues in Application Performance and Application Security
        \item Learnt the distinctive features of Mobile Apps
    \end{itemize}
    \vspace{2em}
    \begin{center}
    \small \textit{Slides used in this presentation are borrowed from \url{http://db-book.com/} with kind permission of the authors.} \\
    \small \textit{Edited and new slides are marked with 'PPD'.}
    \end{center}
\end{frame}

\end{document}

